\documentclass[11pt, letterpaper]{article}

\usepackage[margin=1in]{geometry}
\usepackage{booktabs}
\usepackage{graphicx}
\usepackage{hyperref}
\usepackage{amsmath, amssymb, amsthm}
\usepackage{enumitem}
\usepackage{xcolor}
\usepackage{titlesec}
\usepackage{fancyhdr}
\usepackage{multirow}
\usepackage{caption}
\usepackage{algorithm}
\usepackage{algpseudocode}
\usepackage{listings}
\usepackage{tcolorbox}
\usepackage{float}

% Colors
\definecolor{yellow_conn}{HTML}{F9DF6D}
\definecolor{green_conn}{HTML}{A0C35A}
\definecolor{blue_conn}{HTML}{B0C4EF}
\definecolor{purple_conn}{HTML}{BA81C5}
\definecolor{codebg}{HTML}{F5F5F0}
\definecolor{noteblue}{HTML}{E8F0FE}

\hypersetup{colorlinks=true, linkcolor=blue!60!black, urlcolor=blue!60!black}

\pagestyle{fancy}
\fancyhf{}
\rhead{STA 561D}
\lhead{Infinite Connections --- Study Guide}
\rfoot{\thepage}
\renewcommand{\headrulewidth}{0.4pt}

\titleformat{\section}{\Large\bfseries}{\thesection}{1em}{}
\titleformat{\subsection}{\large\bfseries}{\thesubsection}{1em}{}
\titleformat{\subsubsection}{\normalsize\bfseries}{\thesubsubsection}{1em}{}

\lstset{
    basicstyle=\small\ttfamily,
    backgroundcolor=\color{codebg},
    frame=single,
    framerule=0pt,
    breaklines=true,
    columns=fullflexible,
}

\newtcolorbox{keyidea}{
    colback=noteblue, colframe=blue!40!black,
    fonttitle=\bfseries, title=Key Idea, boxrule=0.5pt
}

\newtcolorbox{formula}{
    colback=yellow_conn!20, colframe=orange!50!black,
    fonttitle=\bfseries, title=Formula, boxrule=0.5pt
}

\newtheorem{definition}{Definition}

\title{\textbf{Infinite Connections}\\[6pt]
\Large Comprehensive Study Guide\\[4pt]
\large AI-Generated NYT Connections Puzzles}
\author{STA 561D --- Duke University}
\date{Spring 2026}

\begin{document}
\maketitle
\thispagestyle{fancy}
\tableofcontents
\newpage

% =========================================================================
\section{The Connections Puzzle}
% =========================================================================

\subsection{What Is It?}

NYT Connections is a daily word puzzle published by \emph{The New York Times}. The player is given \textbf{16 words} and must sort them into \textbf{4 groups of 4}, where each group shares a hidden connection. Groups are color-coded by difficulty:

\begin{center}
\begin{tabular}{@{}lll@{}}
\toprule
Color & Difficulty & Example \\
\midrule
\colorbox{yellow_conn}{Yellow} & Easiest & WAYS TO SAY HELLO: Hi, Hey, Howdy, Yo \\
\colorbox{green_conn}{Green} & Medium & PLANETS: Mars, Venus, Saturn, Jupiter \\
\colorbox{blue_conn}{Blue} & Hard & FIRE \_\_\_: Work, Fly, Place, Side \\
\colorbox{purple_conn}{Purple} & Hardest & TAYLOR SWIFT ALBUMS: Folklore, Reputation, Midnights, Fearless \\
\bottomrule
\end{tabular}
\end{center}

Players get 4 mistakes before losing. A guess of 3 correct out of 4 triggers ``One away!'' feedback.

\subsection{Why Is Generation Hard?}

\begin{enumerate}[nosep]
    \item \textbf{Unique solution.} There are $\binom{16}{4} \times \binom{12}{4} \times \binom{8}{4} = 2{,}627{,}625$ ways to partition 16 words into groups of 4. Exactly \emph{one} must be correct.
    \item \textbf{Difficulty gradient.} Groups must range from obvious to subtle.
    \item \textbf{Deception.} Good puzzles contain plausible-but-wrong groupings (``false groups'') that mislead solvers.
    \item \textbf{Diversity.} Categories should span synonyms, wordplay, fill-in-the-blank, pop culture, etc.
\end{enumerate}

\subsection{The NYT Dataset}

We work with a ground truth dataset of \textbf{554 NYT Connections puzzles} (June 2023 -- December 2024).

\begin{table}[H]
\centering
\caption{NYT dataset summary.}
\begin{tabular}{@{}lr@{}}
\toprule
Metric & Value \\
\midrule
Total puzzles & 554 \\
Unique words & 4,918 \\
Unique categories & 2,054 \\
Multi-word entries & 112 \\
Puzzles with difficulty ratings & 227 \\
Difficulty mean $\pm$ std & $2.94 \pm 0.64$ \\
Most frequent word & BALL (15 appearances) \\
\bottomrule
\end{tabular}
\end{table}

Category type distribution across all 2,216 category slots:
\begin{itemize}[nosep]
    \item Common property: 1,013 (45.7\%) --- e.g., ``THINGS WITH WHEELS''
    \item Synonyms / Slang: 846 (38.2\%) --- e.g., ``WAYS TO SAY GOODBYE''
    \item Fill-in-the-blank: 202 (9.1\%) --- e.g., ``FIRE \_\_\_''
    \item Pop culture: 125 (5.6\%) --- e.g., ``BEYONCE SONGS''
    \item Wordplay / Hidden pattern: 30 (1.4\%) --- e.g., ``B-\_\_\_''
\end{itemize}

% =========================================================================
\section{Foundation: Sentence Embeddings and Cosine Similarity}
% =========================================================================

Every component in this project relies on \textbf{sentence embeddings} --- dense vector representations of words that capture semantic meaning.

\subsection{The MPNET Model}

We use \texttt{all-mpnet-base-v2} from the \texttt{sentence-transformers} library. It maps each word or short phrase to a \textbf{768-dimensional vector}:
\[
    \text{embed}: \text{word} \rightarrow \mathbb{R}^{768}
\]

Words with similar meanings have vectors that point in similar directions. ``PIANO'' and ``GUITAR'' are close in embedding space; ``PIANO'' and ``TRUCK'' are far apart.

\subsection{Cosine Similarity}

\begin{formula}
Given two embedding vectors $\mathbf{a}, \mathbf{b} \in \mathbb{R}^{768}$:
\[
    \cos(\mathbf{a}, \mathbf{b}) = \frac{\mathbf{a} \cdot \mathbf{b}}{\|\mathbf{a}\| \, \|\mathbf{b}\|}
\]
Range: $[-1, 1]$, where $1$ = identical direction, $0$ = orthogonal, $-1$ = opposite.
\end{formula}

In practice, MPNET embeddings are normalized, so cosine similarity reduces to the dot product.

\subsection{Average Pairwise Similarity}

For a group of 4 words with embeddings $\{\mathbf{e}_1, \mathbf{e}_2, \mathbf{e}_3, \mathbf{e}_4\}$, there are $\binom{4}{2} = 6$ pairs:
\begin{formula}
\[
    \bar{s} = \frac{1}{6} \sum_{i < j} \cos(\mathbf{e}_i, \mathbf{e}_j)
\]
\end{formula}

This is the simplest measure of group cohesion. Higher $\bar{s}$ = words are more related.

\begin{keyidea}
Average pairwise cosine similarity is the fundamental building block. It is used for:
\begin{itemize}[nosep]
    \item Selecting the best 4 words from an 8-word candidate pool
    \item Scoring groups in the embedding solver
    \item Assigning difficulty colors
    \item Evaluating puzzle quality
\end{itemize}
\end{keyidea}

% =========================================================================
\section{Puzzle Generation Pipeline}
% =========================================================================

The pipeline creates puzzles through six stages. The LLM provides creativity; MPNET provides consistency.

\subsection{Stage 1: Group Creation (LLM + Embedding Selection)}

Each puzzle is built \textbf{one group at a time} across 4 LLM calls.

\begin{algorithm}[H]
\caption{Iterative Group Creation}
\begin{algorithmic}[1]
\State $\text{groups} \gets []$
\For{$i = 1$ to $4$}
    \State $\text{seeds} \gets \text{sample 4 random words from NYT word bank}$
    \State \textbf{LLM call}: generate category name + 8 candidate words (story-injected)
    \State $\text{pool} \gets 8$ candidate words
    \State $\text{best4}, \text{score} \gets \textsc{SelectBest}(\text{pool}, n=4)$ \Comment{MPNET}
    \State Remove any words already in previous groups
    \State Append $(\text{category}, \text{best4}, \text{score})$ to groups
\EndFor
\State \Return groups
\end{algorithmic}
\end{algorithm}

\subsubsection{Embedding Selection: 8 $\rightarrow$ 4}

The LLM proposes 8 candidates; we use MPNET to select the 4 most internally similar:

\begin{algorithm}[H]
\caption{SelectBest: Choose $k$ most cohesive words from $n$}
\begin{algorithmic}[1]
\Require Words $W = \{w_1, \ldots, w_n\}$, target size $k$
\State Encode all words: $\mathbf{E} = \{\text{embed}(w_i)\}$
\State $\text{best\_score} \gets -1$
\For{each $C \in \binom{W}{k}$} \Comment{$\binom{8}{4} = 70$ combinations}
    \State $s \gets$ average pairwise cosine similarity of $C$
    \If{$s > \text{best\_score}$}
        \State $\text{best\_score} \gets s$; $\text{best\_combo} \gets C$
    \EndIf
\EndFor
\State \Return best\_combo, best\_score
\end{algorithmic}
\end{algorithm}

\begin{keyidea}
The LLM is not allowed to pick the final 4 words. LLMs are bad at estimating semantic similarity and make inconsistent selections. MPNET selection is deterministic and correlates with human difficulty perception.
\end{keyidea}

\subsubsection{Story Injection for Diversity}

Without intervention, the LLM generates the same $\sim$20 categories repeatedly (BOARD GAMES, FRUITS, COLORS, \ldots). The \textbf{story injection} technique solves this:

\begin{quote}
\emph{``First, write a short story (2--3 sentences) using these words: [4 random words from the NYT word bank]. Then use that story as inspiration for creating your category.''}
\end{quote}

The random seed words force the LLM into different creative directions each time.

\subsection{Stage 2: False-Group Method (Primary)}

This produces the highest-quality puzzles. A user study found that false-group puzzles beat NYT originals in user preference 42.86\% of the time.

\begin{algorithm}[H]
\caption{False-Group Pipeline}
\begin{algorithmic}[1]
\State \textbf{LLM call}: Generate root group $R = (c_R, \{w_1, w_2, w_3, w_4\})$
\State \quad Each $w_i$ must have an alternate meaning $m_i$
\State $\text{groups} \gets []$
\For{each $(w_i, m_i) \in R$}
    \State \textbf{LLM call}: Generate category inspired by meaning $m_i$, with 8 candidates
    \State $\text{best4} \gets \textsc{SelectBest}(\text{candidates}, 4)$ \Comment{MPNET}
    \State Append to groups
\EndFor
\State \Return groups \Comment{Root group $R$ is the decoy --- discarded}
\end{algorithmic}
\end{algorithm}

\textbf{Example:} Root group ``THINGS THAT ARE ROUND'' $\rightarrow$ \{BALL, GLOBE, RING, WHEEL\}.
\begin{itemize}[nosep]
    \item BALL $\xrightarrow{\text{alternate}}$ ``a formal dance event'' $\rightarrow$ generates DANCE category
    \item GLOBE $\xrightarrow{\text{alternate}}$ ``a theater or news org'' $\rightarrow$ generates THEATER category
    \item RING $\xrightarrow{\text{alternate}}$ ``the sound a phone makes'' $\rightarrow$ generates SOUNDS category
    \item WHEEL $\xrightarrow{\text{alternate}}$ ``to turn or rotate'' $\rightarrow$ generates MOVEMENT category
\end{itemize}

The solver sees 16 words that \emph{look like} they could form a ``ROUND THINGS'' group, but that group doesn't exist. This creates natural deception.

\subsection{Stage 3: Editor Pass}

A second LLM call reviews the complete puzzle:
\begin{itemize}[nosep]
    \item Checks if each category name accurately describes its 4 words
    \item Rewrites inaccurate names
    \item Flags groups where a word doesn't fit
\end{itemize}

\subsection{Stage 4: Difficulty Color Assignment}

\begin{formula}
For each group, compute average pairwise cosine similarity $\bar{s}$. Sort groups by $\bar{s}$ descending:
\[
    \text{highest } \bar{s} \rightarrow \text{Yellow},\quad
    \text{next} \rightarrow \text{Green},\quad
    \text{next} \rightarrow \text{Blue},\quad
    \text{lowest } \bar{s} \rightarrow \text{Purple}
\]
\end{formula}

\textbf{Intuition:} High semantic similarity = words obviously relate = easy. Low similarity = the connection is subtle or relies on wordplay = hard.

Reference thresholds from the literature (validated against NYT data):

\begin{table}[H]
\centering
\caption{Color thresholds: paper values vs.\ our empirical measurements on 554 NYT puzzles.}
\begin{tabular}{@{}lcccc@{}}
\toprule
& \multicolumn{2}{c}{Paper (expected)} & \multicolumn{2}{c}{Our measurement} \\
\cmidrule(lr){2-3} \cmidrule(lr){4-5}
Color & Mean $\bar{s}$ & Variance & Mean $\bar{s}$ & Std \\
\midrule
\colorbox{yellow_conn}{Yellow} & 0.40 & 0.0285 & 0.389 & 0.089 \\
\colorbox{green_conn}{Green}   & 0.35 & 0.0214 & 0.315 & 0.060 \\
\colorbox{blue_conn}{Blue}     & 0.29 & 0.0123 & 0.268 & 0.050 \\
\colorbox{purple_conn}{Purple} & 0.27 & 0.0108 & 0.225 & 0.049 \\
\bottomrule
\end{tabular}
\end{table}

NYT's color ordering matches pure similarity ranking only \textbf{43.9\%} of the time, confirming that difficulty has a knowledge/wordplay component that embeddings alone cannot capture.

\subsection{Stage 5: Deduplication}

Each generated puzzle's 16-word set is compared against all 554 known NYT puzzles:
\[
    \text{overlap}(P_\text{gen}, P_\text{nyt}) = |W_\text{gen} \cap W_\text{nyt}|
\]
If $\text{overlap} > 6$ for any NYT puzzle, the generated puzzle is flagged as a duplicate.

\subsection{Pipeline Summary}

\begin{center}
\begin{tabular}{@{}lll@{}}
\toprule
Stage & Who does it & Why \\
\midrule
Group creation & LLM (creative) & Invents categories and candidate words \\
Word selection & MPNET (deterministic) & Picks the most cohesive 4 from 8 \\
Editor pass & LLM (quality check) & Fixes inaccurate category names \\
Color assignment & MPNET (scoring) & Maps similarity to difficulty \\
Deduplication & Set intersection & Prevents copying NYT puzzles \\
\bottomrule
\end{tabular}
\end{center}

% =========================================================================
\section{Solver Algorithms}
% =========================================================================

Three independent solvers verify that each puzzle has exactly one valid solution.

\subsection{Solver 1: Embedding Solver (Greedy)}

\textbf{Source:} ``Missed Connections'' paper (arXiv:2404.11730).

\begin{algorithm}[H]
\caption{Embedding Solver}
\begin{algorithmic}[1]
\Require 16 words $W$
\State Encode all words: $\mathbf{E} = \{\text{embed}(w_i)\}_{i=1}^{16}$
\State Compute $16 \times 16$ cosine similarity matrix $S$
\State Score all $\binom{16}{4} = 1{,}820$ possible groups by $\bar{s}$
\State Sort groups by score (descending)
\State $\text{selected} \gets []$; $\text{used} \gets \emptyset$
\For{each $(\text{group}, \text{score})$ in sorted order}
    \If{group has no overlap with used}
        \State Append group to selected; add words to used
    \EndIf
    \If{$|\text{selected}| = 4$} \textbf{break} \EndIf
\EndFor
\State \Return selected
\end{algorithmic}
\end{algorithm}

\textbf{Complexity:} $O\bigl(\binom{16}{4}\bigr) = O(1{,}820)$ group evaluations + sorting.\\
\textbf{Runtime:} $\sim$0.019s per puzzle.\\
\textbf{Accuracy:} 2.0\% full-puzzle solve rate on 554 NYT puzzles.

\textbf{Limitation:} Greedy selection can ``use up'' key words in a wrong group, causing cascading errors. The solver never backtracks.

\subsection{Solver 2: Clustering Solver (Beam Search)}

\textbf{Source:} ``Deceptively Simple'' paper (arXiv:2412.01621).

This solver uses a more sophisticated scoring function and beam search to avoid greedy errors.

\subsubsection{Group Similarity Score ($G$)}

\begin{formula}
For a group of 4 word embeddings $\{\mathbf{e}_1, \ldots, \mathbf{e}_4\}$:
\[
    G = 0.4 \cdot I + 0.3 \cdot s + 0.3 \cdot V
\]
where:
\begin{align}
    I &= -\sum_{i=1}^{4} \|\mathbf{e}_i - \boldsymbol{\mu}\|^2
    & \text{(negative inertia, } \boldsymbol{\mu} = \text{centroid)} \\[4pt]
    s &= \min_{i < j} \cos(\mathbf{e}_i, \mathbf{e}_j)
    & \text{(minimum pairwise similarity)} \\[4pt]
    V &= \frac{\text{mean}(P)}{1 + \text{var}(P)}
    & \text{(stability-weighted mean)}
\end{align}
$P = \{\cos(\mathbf{e}_i, \mathbf{e}_j) : i < j\}$ is the set of all 6 pairwise similarities.
\end{formula}

\textbf{Intuition for each component:}
\begin{itemize}[nosep]
    \item $I$ (inertia): How tightly clustered are the words? Higher $|I|$ means tighter clustering. The negative sign makes ``tighter = higher score.''
    \item $s$ (min similarity): Penalizes groups where even one pair is weakly related. A group is only as strong as its weakest link.
    \item $V$ (stability): Rewards consistent similarity across all pairs. A group where all pairs have similarity 0.3 scores higher than one with similarities \{0.5, 0.5, 0.5, 0.1, 0.1, 0.1\}, even if the means are equal.
\end{itemize}

\subsubsection{Penalty Score ($P$)}

\begin{formula}
After selecting a candidate group $C$ with centroid $\boldsymbol{\mu}_C$, compute how strongly the remaining words $R$ are attracted to it:
\[
    P = \frac{1}{|R|} \sum_{r \in R} \cos(\boldsymbol{\mu}_C, \mathbf{r})
\]
Lower penalty = better separation between this group and everything else.
\end{formula}

\subsubsection{Beam Search}

\begin{algorithm}[H]
\caption{Clustering Solver (Beam Search)}
\begin{algorithmic}[1]
\Require 16 words $W$, beam width $b = 10$
\State $\text{beam} \gets [(\text{groups}=[], \text{remaining}=W, \text{score}=0)]$
\For{step $= 1$ to $4$}
    \State $\text{candidates} \gets []$
    \For{each state $(\text{groups}, \text{remaining}, \text{cumul})$ in beam}
        \For{each group $C \in \binom{\text{remaining}}{4}$}
            \State $g \gets G(C)$ \Comment{Group Similarity Score}
            \State $p \gets P(C, \text{remaining} \setminus C)$ \Comment{Penalty Score}
            \State $\text{new\_score} \gets \text{cumul} + g - p$
            \State Add $(\text{groups} \cup \{C\},\ \text{remaining} \setminus C,\ \text{new\_score})$ to candidates
        \EndFor
    \EndFor
    \State Sort candidates by score (descending)
    \State $\text{beam} \gets$ top $b$ candidates
\EndFor
\State \Return best complete solution from beam
\end{algorithmic}
\end{algorithm}

\textbf{Complexity:} At each step, up to $b \times \binom{|\text{remaining}|}{4}$ evaluations. Step 1: $10 \times 1{,}820$; Step 2: $10 \times 495$; Step 3: $10 \times 70$; Step 4: $10 \times 1$.\\
\textbf{Runtime:} $\sim$0.95s per puzzle ($\sim$50$\times$ slower than embedding solver).\\
\textbf{Accuracy:} 3.4\% full-puzzle solve rate.

\subsection{Solver 3: LLM Solver (Chain-of-Thought)}

\textbf{Source:} ``Missed Connections'' paper.

Presents all 16 words to Claude with a chain-of-thought prompt:
\begin{enumerate}[nosep]
    \item Look for obvious groupings first
    \item Consider alternate meanings and wordplay
    \item Every word belongs to exactly one group
    \item Each group has exactly 4 words
\end{enumerate}

This is the most ``human-like'' solver but expensive (requires an API call per puzzle). Used as a tiebreaker only.

\subsection{Solver Performance Comparison}

\begin{table}[H]
\centering
\caption{Benchmark on 554 NYT puzzles.}
\begin{tabular}{@{}lcc@{}}
\toprule
& Embedding & Clustering \\
\midrule
Full solve rate & 2.0\% (11/554) & 3.4\% (19/554) \\
Avg groups correct & 0.57 / 4 & 0.71 / 4 \\
\midrule
\colorbox{yellow_conn}{Yellow} accuracy & 24.5\% & 28.0\% \\
\colorbox{green_conn}{Green} accuracy & 17.7\% & 22.6\% \\
\colorbox{blue_conn}{Blue} accuracy & 9.8\% & 13.0\% \\
\colorbox{purple_conn}{Purple} accuracy & 5.1\% & 7.6\% \\
\midrule
Runtime (554 puzzles) & 10.6\,s & 525\,s \\
\bottomrule
\end{tabular}
\end{table}

Both solvers show the expected difficulty gradient. The clustering solver outperforms across all colors at $\sim$50$\times$ the compute cost. They agreed on a full solve for only 7 puzzles (1.3\%), confirming they find complementary solutions.

% =========================================================================
\section{Roundtable Validation}
% =========================================================================

\begin{keyidea}
This is the project's most original contribution. The papers each address generation OR solving; we combine them by using solver agreement as a quality gate for generated puzzles.
\end{keyidea}

\subsection{The Core Idea}

A valid Connections puzzle must have \textbf{exactly one correct solution}. If two independent solvers, using completely different algorithms, both find the same grouping --- and that grouping matches the intended answer --- then the puzzle is unambiguous.

If they disagree, the puzzle has competing interpretations and should be discarded.

\subsection{Validation Logic}

\begin{algorithm}[H]
\caption{Roundtable Validation}
\begin{algorithmic}[1]
\Require Puzzle $P$ with intended groups $A$
\State $S_1 \gets \textsc{EmbeddingSolver}(P.\text{words})$
\State $S_2 \gets \textsc{ClusteringSolver}(P.\text{words})$
\State Normalize $S_1, S_2$ to sets of frozensets (order-independent)
\State $\text{agree} \gets (S_1 = S_2)$ \Comment{Same 4 groups?}
\State $c_1 \gets |\{g \in S_1 : g \in A\}|$ \Comment{How many match the answer?}
\State $c_2 \gets |\{g \in S_2 : g \in A\}|$
\State $\text{valid} \gets \text{agree} \textbf{ and } (c_1 = 4 \textbf{ or } c_2 = 4)$
\State \Return valid, agree, $c_1$, $c_2$
\end{algorithmic}
\end{algorithm}

\subsection{Why Two Solvers?}

The embedding solver and clustering solver use different algorithms:
\begin{itemize}[nosep]
    \item Embedding: greedy, optimizes average pairwise similarity
    \item Clustering: beam search, optimizes $G - P$ composite score
\end{itemize}

They have different failure modes. When both converge to the same answer, it is strong evidence that the answer is unique.

% =========================================================================
\section{Evaluation Metrics}
% =========================================================================

\subsection{Group Similarity Score ($G$)}

Identical to the clustering solver's scoring function (Section 4.2.1). Used as an independent quality metric for generated groups.

\subsection{Penalty Score ($P$)}

Identical to Section 4.2.2. Measures group--group separation.

\subsection{Puzzle-Level Quality}

For a complete puzzle:
\begin{itemize}[nosep]
    \item \textbf{Overall avg similarity:} mean of $\bar{s}$ across all 4 groups
    \item \textbf{Similarity spread:} $\max(\bar{s}) - \min(\bar{s})$ across groups. Should be positive (yellow $>$ purple).
    \item \textbf{Solver agreement rate:} fraction of generated puzzles where solvers converge
    \item \textbf{Category diversity:} number of unique category names across the dataset
\end{itemize}

% =========================================================================
\section{System Architecture}
% =========================================================================

\subsection{Data Flow}

\begin{center}
\begin{tabular}{@{}c@{}}
\texttt{LLM (Claude)} $\xrightarrow{\text{8 candidates}}$
\texttt{MPNET (select 4)} $\xrightarrow{\text{groups}}$
\texttt{Editor (LLM)} $\xrightarrow{\text{reviewed}}$
\texttt{Difficulty (MPNET)} \\[8pt]
$\downarrow$ \\[4pt]
\texttt{Roundtable Validator} $\xrightarrow{\text{agree?}}$
$\begin{cases} \text{Yes} \rightarrow \text{Accept puzzle} \\ \text{No} \rightarrow \text{Discard} \end{cases}$
\end{tabular}
\end{center}

\subsection{Dry-Run Mode}

All LLM calls go through \texttt{LLMClient}, which checks:
\begin{lstlisting}
if self.dry_run:
    return MockLLMClient.respond(system, user)  # no API call
\end{lstlisting}

\texttt{MockLLMClient} detects the request type from prompt keywords and returns realistic hardcoded responses. MPNET always runs for real --- only the LLM is mocked.

\begin{table}[H]
\centering
\caption{What is real vs.\ mocked in dry-run mode.}
\begin{tabular}{@{}lcc@{}}
\toprule
Component & Dry-run & Real mode \\
\midrule
Category ideas + word pools & Mock & Claude API \\
Word selection (8 $\to$ 4) & Real MPNET & Real MPNET \\
Editor pass & Mock & Claude API \\
Difficulty colors & Real MPNET & Real MPNET \\
Embedding solver & Real MPNET & Real MPNET \\
Clustering solver & Real MPNET & Real MPNET \\
LLM solver & Mock & Claude API \\
\bottomrule
\end{tabular}
\end{table}

\subsection{File Map}

\begin{table}[H]
\centering
\small
\caption{Key source files and their responsibilities.}
\begin{tabular}{@{}lp{8cm}@{}}
\toprule
File & Role \\
\midrule
\texttt{src/config.py} & DRY\_RUN flag, paths, thresholds, model names \\
\texttt{src/llm\_client.py} & LLM abstraction + MockLLMClient \\
\texttt{src/generator/pipeline.py} & Orchestrates all 6 generation stages \\
\texttt{src/generator/group\_creator.py} & LLM calls + MPNET embedding selection \\
\texttt{src/generator/prompts.py} & All prompt templates \\
\texttt{src/generator/difficulty.py} & Color assignment via cosine similarity \\
\texttt{src/generator/deduplicator.py} & Overlap check vs.\ NYT puzzles \\
\texttt{src/solvers/embedding\_solver.py} & Greedy solver \\
\texttt{src/solvers/clustering\_solver.py} & Beam search solver \\
\texttt{src/solvers/roundtable.py} & Multi-solver convergence validator \\
\texttt{src/evaluation/metrics.py} & $G$, $P$, and quality scores \\
\texttt{tests/test\_pipeline.py} & 31 tests covering all modules \\
\bottomrule
\end{tabular}
\end{table}

% =========================================================================
\section{Key Formulas Reference Card}
% =========================================================================

\begin{tcolorbox}[colback=codebg, colframe=black!50, title=\textbf{All Formulas in One Place}]

\textbf{Cosine similarity:}
$\displaystyle \cos(\mathbf{a}, \mathbf{b}) = \frac{\mathbf{a} \cdot \mathbf{b}}{\|\mathbf{a}\| \, \|\mathbf{b}\|}$
\vspace{6pt}

\textbf{Average pairwise similarity (4 words):}
$\displaystyle \bar{s} = \frac{1}{6}\sum_{i<j} \cos(\mathbf{e}_i, \mathbf{e}_j)$
\vspace{6pt}

\textbf{Group Similarity Score:}
$\displaystyle G = 0.4 \cdot \underbrace{\bigl(-\textstyle\sum_i \|\mathbf{e}_i - \boldsymbol{\mu}\|^2\bigr)}_{I\text{ (neg.\ inertia)}} + 0.3 \cdot \underbrace{\min_{i<j}\cos(\mathbf{e}_i,\mathbf{e}_j)}_{s\text{ (weakest link)}} + 0.3 \cdot \underbrace{\frac{\text{mean}(P)}{1+\text{var}(P)}}_{V\text{ (stability)}}$
\vspace{6pt}

\textbf{Penalty Score:}
$\displaystyle P = \frac{1}{|R|}\sum_{r \in R} \cos(\boldsymbol{\mu}_C, \mathbf{r})$
\vspace{6pt}

\textbf{Beam search objective:} maximize $\displaystyle\sum_{\text{groups}} (G - P)$
\vspace{6pt}

\textbf{Color assignment:} sort groups by $\bar{s}$ descending $\to$ Yellow, Green, Blue, Purple
\vspace{6pt}

\textbf{Validation:} $\text{valid} = (\text{solver}_1 \equiv \text{solver}_2) \wedge (\exists\, \text{solver with 4/4 correct})$

\end{tcolorbox}

% =========================================================================
\section{What Is Original vs.\ From the Literature}
% =========================================================================

\begin{table}[H]
\centering
\small
\caption{Provenance of each component.}
\begin{tabular}{@{}lll@{}}
\toprule
Component & Status & Source of ideas \\
\midrule
Embedding solver & Reimplementation & ``Missed Connections'' (2404.11730) \\
Clustering solver ($G$-$P$ + beam) & Reimplementation & ``Deceptively Simple'' (2412.01621) \\
False-group pipeline & Reimplementation & ``Making New Connections'' (2407.11240) \\
Story injection & Reimplementation & ``Making New Connections'' \\
Color thresholds & Constants from paper & ``Making New Connections'' Table 1 \\
LLM solver (CoT) & Reimplementation & ``Missed Connections'' \\
\midrule
Roundtable validator & \textbf{Novel} & Original to this project \\
Prompt templates & \textbf{Novel} & Original text \\
MockLLMClient / dry-run & \textbf{Novel} & Original engineering \\
Web app & \textbf{Novel} & Original code \\
End-to-end integration & \textbf{Novel} & Papers do generation OR solving; we combine them \\
\bottomrule
\end{tabular}
\end{table}

All code is written from scratch. No code was copied from any reference repository.

% =========================================================================
\section{References}
% =========================================================================

\begin{enumerate}[nosep, label={[\arabic*]}]
    \item A.\ Conneau et al., ``Making New Connections: LLM-based puzzle generation,'' arXiv:2407.11240, 2024. \emph{(Generation pipeline, false groups, story injection, color thresholds.)}
    \item M.\ Todd et al., ``Missed Connections: Lateral thinking in LLMs,'' arXiv:2404.11730, 2024. \emph{(Embedding solver, LLM solver baselines, CoT prompts.)}
    \item R.\ Li et al., ``Deceptively Simple: A scoring and search approach,'' arXiv:2412.01621, 2024. \emph{($G = 0.4I + 0.3s + 0.3V$ formula, penalty score, beam search.)}
    \item J.\ Patel et al., ``Connecting the Dots: A knowledge taxonomy,'' arXiv:2406.11012, 2024. \emph{(Category types, human evaluation methodology.)}
\end{enumerate}

\end{document}
